% Vector addition and dot product on the Kites RISC-V simulator: C code, scalar and vector
% assembly, performance comparison, and the implementation of the vector memory instructions.
% Self-contained: compile with  pdflatex vector_add_dot_report.tex
% All numbers were measured with kites-cli on the single-cycle VM (benchmarks/results/).
\documentclass[11pt,a4paper]{article}

\usepackage[margin=2.2cm]{geometry}
\usepackage[T1]{fontenc}
\usepackage{lmodern}
\usepackage{amsmath}
\usepackage{booktabs}
\usepackage{array}
\usepackage{xcolor}
\usepackage{listings}
\usepackage{caption}
\usepackage{float}
\usepackage[hidelinks]{hyperref}

\newcommand{\code}[1]{\texttt{#1}}
\captionsetup{font=small,labelfont=bf}

\definecolor{cmt}{RGB}{110,110,104}
\definecolor{kw}{RGB}{30,90,170}
\lstdefinestyle{asm}{
  basicstyle=\ttfamily\footnotesize,
  commentstyle=\color{cmt}\itshape,
  morecomment=[l]{\#},
  alsoletter={.},
  keywordstyle=\color{kw}\bfseries,
  morekeywords={vsetvli,vle32.v,vse32.v,vadd.vv,vmul.vv,vredsum.vs,vmv.s.x,vmv.x.s},
  columns=fullflexible,
  keepspaces=true,
  frame=single,
  rulecolor=\color{black!25},
  xleftmargin=1em,
  framexleftmargin=0.8em
}
\lstdefinestyle{c}{
  language=C,
  basicstyle=\ttfamily\footnotesize,
  keywordstyle=\color{kw}\bfseries,
  commentstyle=\color{cmt}\itshape,
  columns=fullflexible,
  frame=single,
  rulecolor=\color{black!25},
  xleftmargin=1em,
  framexleftmargin=0.8em
}

\title{Vector Addition and Dot Product on Kites:\\ Scalar vs.\ Vector Instructions}
\author{RISC-V Vector Extension (RVV 1.0) project}
\date{\today}

\begin{document}
\maketitle

%==============================================================================================
\section{Introduction}
%==============================================================================================

This report compares two common operations, \emph{vector addition} and the \emph{dot product},
written once with the regular RISC-V integer instructions (RV64IM, \emph{scalar}) and once
with the RISC-V Vector extension (RVV~1.0, \emph{vector}). Both versions of each operation read
the same 32-bit integer arrays and write their result to the same address, and both were run on
the single-cycle VM of the Kites simulator, which executes one instruction per step.

The vector register length is VLEN~=~128~bits. With 32-bit elements (SEW~=~32) and groups of 8
vector registers (LMUL~=~8), one vector instruction processes up to
$\mathrm{VLMAX} = \mathrm{LMUL}\cdot\mathrm{VLEN}/\mathrm{SEW} = 8\cdot128/32 = 32$ elements.

%==============================================================================================
\section{Vector addition}
%==============================================================================================

\subsection{C code}

\begin{lstlisting}[style=c]
int32_t a[N], b[N], c[N];
for (int i = 0; i < N; i++)
    c[i] = a[i] + b[i];
\end{lstlisting}

\subsection{Mapping to assembly}

\begin{minipage}[t]{0.47\textwidth}
\textbf{Scalar (RV64IM)}
\begin{lstlisting}[style=asm]
    la   a0, a          # pointer to a[i]
    la   a1, b          # pointer to b[i]
    la   a2, c          # pointer to c[i]
    li   a3, N          # elements left
loop:
    lw   t0, 0(a0)      # load a[i]
    lw   t1, 0(a1)      # load b[i]
    add  t0, t0, t1     # a[i] + b[i]
    sw   t0, 0(a2)      # store c[i]
    addi a0, a0, 4      # next element
    addi a1, a1, 4
    addi a2, a2, 4
    addi a3, a3, -1     # one element done
    bnez a3, loop
\end{lstlisting}
\end{minipage}\hfill
\begin{minipage}[t]{0.50\textwidth}
\textbf{Vector (RVV 1.0)}
\begin{lstlisting}[style=asm]
    la   a0, a
    la   a1, b
    la   a2, c
    li   a3, N
loop:
    vsetvli t0, a3, e32, m8, ta, ma # vl
    vle32.v v0, (a0)      # load vl a's
    vle32.v v8, (a1)      # load vl b's
    vadd.vv v16, v0, v8   # vl additions
    vse32.v v16, (a2)     # store vl c's
    slli t1, t0, 2        # vl * 4 bytes
    add  a0, a0, t1       # advance pointers
    add  a1, a1, t1
    add  a2, a2, t1
    sub  a3, a3, t0       # vl elements done
    bnez a3, loop
\end{lstlisting}
\end{minipage}

\begin{itemize}
  \item \textbf{Scalar mapping.} Each iteration of the C loop becomes one iteration of the
        assembly loop: two loads read \code{a[i]} and \code{b[i]}, one \code{add} computes the sum,
        and one store writes \code{c[i]}. The index \code{i} is replaced by three pointers that
        advance by 4 bytes, and the loop condition becomes a counter that is decremented and tested
        by \code{bnez}. Each element therefore costs 9 instructions, of which only one is the
        addition.
  \item \textbf{Vector mapping.} The loop is \emph{strip-mined}: in each pass \code{vsetvli}
        sets the vector length $vl = \min(\text{elements left}, 32)$. The two \code{vle32.v}
        instructions load $vl$ elements of \code{a} and \code{b}, a single \code{vadd.vv} adds all
        of them, and \code{vse32.v} stores the $vl$ results. The pointers advance by $vl \cdot 4$
        bytes and the counter by $vl$. One pass of 11 instructions handles up to 32 elements, and
        the last pass automatically handles any remainder.
\end{itemize}

%==============================================================================================
\section{Dot product}
%==============================================================================================

\subsection{C code}

\begin{lstlisting}[style=c]
int32_t a[N], b[N];
int32_t sum = 0;
for (int i = 0; i < N; i++)
    sum += a[i] * b[i];
\end{lstlisting}

\subsection{Mapping to assembly}

\begin{minipage}[t]{0.47\textwidth}
\textbf{Scalar (RV64IM)}
\begin{lstlisting}[style=asm]
    la   a0, a
    la   a1, b
    la   a2, out
    li   a3, N
    li   a4, 0          # sum = 0
loop:
    lw   t0, 0(a0)      # a[i]
    lw   t1, 0(a1)      # b[i]
    mul  t0, t0, t1     # a[i] * b[i]
    add  a4, a4, t0     # sum += product
    addi a0, a0, 4
    addi a1, a1, 4
    addi a3, a3, -1
    bnez a3, loop
    sw   a4, 0(a2)      # out = sum
\end{lstlisting}
\end{minipage}\hfill
\begin{minipage}[t]{0.50\textwidth}
\textbf{Vector (RVV 1.0)}
\begin{lstlisting}[style=asm]
    la   a0, a
    la   a1, b
    la   a2, out
    li   a3, N
    vsetvli t0, x0, e32, m1, ta, ma
    vmv.s.x v24, zero     # v24[0] = 0 (sum)
loop:
    vsetvli t0, a3, e32, m8, ta, ma
    vle32.v v0, (a0)
    vle32.v v8, (a1)
    vmul.vv v16, v0, v8   # vl products
    vredsum.vs v24, v16, v24 # sum += all vl
    slli t1, t0, 2
    add  a0, a0, t1
    add  a1, a1, t1
    sub  a3, a3, t0
    bnez a3, loop
    vmv.x.s a4, v24       # a4 = v24[0]
    sw   a4, 0(a2)        # out = sum
\end{lstlisting}
\end{minipage}

\begin{itemize}
  \item \textbf{Scalar mapping.} The variable \code{sum} lives in register \code{a4}. Each
        iteration loads one pair, multiplies it with \code{mul} and adds the product to \code{a4};
        after the loop a single \code{sw} writes the sum to memory. Each element costs 8
        instructions.
  \item \textbf{Vector mapping.} The running sum is kept in element 0 of \code{v24}, which is
        set to 0 before the loop with \code{vmv.s.x}. Each pass loads $vl$ pairs, multiplies them
        all with one \code{vmul.vv}, and adds all $vl$ products to \code{v24[0]} with one
        reduction instruction, \code{vredsum.vs}. After the loop \code{vmv.x.s} copies the sum to a
        scalar register so that it can be stored. One pass of 10 instructions handles up to 32
        elements.
\end{itemize}

%==============================================================================================
\section{Performance comparison}
%==============================================================================================

\subsection{Metrics}

\begin{enumerate}
  \item \textbf{Dynamic instruction count (IC) and speedup.} The IC is the number of instructions
        actually executed, counted by the simulator. The speedup is
        $S = \mathrm{IC}_{\text{scalar}} / \mathrm{IC}_{\text{vector}}$. On the single-cycle VM every
        instruction takes one step, so this is also the speedup in simulated execution time.
  \item \textbf{Instructions per element (IPE).} $\mathrm{IPE} = \mathrm{IC}/N$, the cost of
        processing one element, independent of the array size.
  \item \textbf{Memory instructions.} The number of load and store instructions executed. Both
        versions move exactly the same number of bytes, so this shows how much data each memory
        instruction carries.
  \item \textbf{Branches executed.} The number of conditional branches, which equals the number of
        loop iterations and measures the control overhead.
\end{enumerate}

From the programs above, the instruction counts are $7 + 9N$ (scalar) and
$7 + 11\lceil N/32\rceil$ (vector) for addition, and $9 + 8N$ and $11 + 10\lceil N/32\rceil$
for the dot product. The measured counts match these formulas exactly. For $N=256$, for example,
addition takes $7 + 9\cdot256 = 2311$ scalar and $7 + 11\cdot8 = 95$ vector instructions.

\subsection{Results}

\begin{table}[H]
\centering\small\setlength{\tabcolsep}{5pt}
\caption{Scalar vs.\ vector (VLEN~=~128). Pairs are scalar / vector.}
\label{tab:results}
\begin{tabular}{lrrrcccc}
\toprule
Operation & $N$ & IC scalar & IC vector & \textbf{Speedup} & IPE & memory instr. & branches \\
\midrule
Vector addition & 256     & 2{,}311  & 95      & \textbf{24.3$\times$} & 9.03 / 0.37 & 768 / 24        & 256 / 8 \\
Vector addition & 4{,}096 & 36{,}871 & 1{,}415 & \textbf{26.1$\times$} & 9.00 / 0.35 & 12{,}288 / 384  & 4{,}096 / 128 \\
\addlinespace
Dot product     & 256     & 2{,}057  & 91      & \textbf{22.6$\times$} & 8.04 / 0.36 & 513 / 17        & 256 / 8 \\
Dot product     & 4{,}096 & 32{,}777 & 1{,}291 & \textbf{25.4$\times$} & 8.00 / 0.32 & 8{,}193 / 257   & 4{,}096 / 128 \\
\bottomrule
\end{tabular}
\end{table}

\begin{itemize}
  \item \textbf{Speedup.} The vector versions execute 22--26 times fewer instructions. For long
        arrays the speedup approaches the ratio of the loop lengths times VLMAX:
        $9\cdot32/11 = 26.2$ for addition and $8\cdot32/10 = 25.6$ for the dot product.
  \item \textbf{Instructions per element.} The scalar loops need 8--9 instructions for every
        element, while the vector loops need only about 0.35, because each vector instruction works
        on 32 elements at once.
  \item \textbf{Memory instructions.} Both versions move the same data (3{,}072 bytes for addition
        at $N=256$), but one \code{vle32.v} loads 128 bytes where an \code{lw} loads 4, so the vector
        versions need about 32 times fewer memory instructions.
  \item \textbf{Branches.} The loop overhead (pointer updates, counter and branch) is paid once per
        32 elements instead of once per element, so the branches fall from 256 to 8.
  \item \textbf{Longer vector registers.} The same vector code needs no change for a larger VLEN.
        At VLEN~=~1024 the addition of 256 elements takes 18 instructions instead of 95
        (128$\times$ faster than scalar), and the dot product takes 21 (98$\times$).
\end{itemize}

%==============================================================================================
\section{Implementation of the vector memory instructions}
%==============================================================================================

\begin{itemize}
  \item \textbf{New files.} All vector loads and stores are implemented in
        \code{vector\_memory.h} and \code{vector\_memory.cpp} (in \code{src/processor/vector/}). The
        assembler, the disassembler and all five processor models share this code.
  \item \textbf{Encoding.} Vector loads and stores reuse the floating-point load/store opcodes; the
        \code{width} field distinguishes them from \code{flw} and \code{fld}.
  \item \textbf{Supported forms.} Unit-stride, strided, indexed (ordered and unordered), segment,
        whole-register, mask and fault-only-first loads and stores, for 8-, 16-, 32- and 64-bit
        elements.
  \item \textbf{Assembler.} A mnemonic table converts each instruction name to its encoding, and
        the parser checks the operands of each form. The disassembler prints the instructions in
        the GNU assembler format.
  \item \textbf{Execution.} \code{ExecuteVectorMemory()} first checks that the instruction is
        legal for the current \code{vtype} (register-group size and alignment, mask register use),
        then accesses the active elements one by one, from \code{vstart} to \code{vl}.
  \item \textbf{Memory path.} Every element goes through the normal memory controller and caches,
        so vector accesses appear in the cache statistics and can be undone like scalar ones.
  \item \textbf{Policies.} Elements past \code{vl} and masked-off elements are left unchanged, as
        in the reference simulator Spike. An invalid address raises an access fault, except in
        fault-only-first loads, where a fault after element 0 shortens \code{vl} instead.
  \item \textbf{Processors.} The whole access is performed in the execute stage, on the
        single-cycle VM and on the pipelines alike, which keeps scalar and vector memory accesses
        in program order.
\end{itemize}

%==============================================================================================
\section{Testing suite for the load and store instructions}
%==============================================================================================

\begin{itemize}
  \item \textbf{Unit tests:} 30 tests in \code{test/vector\_load\_store\_test.cpp} check the
        encodings, the assembler and disassembler, and every addressing mode. All pass.
  \item \textbf{End-to-end tests:} 15 small programs in \code{examples/vector\_ldst\_tests/} were
        run on the single-cycle VM and both hazard-detecting pipelines, and their memory was
        compared with the expected bytes. All 45 runs pass.
  \item \textbf{$vl = 0$:} loads and stores change nothing.
  \item \textbf{Tail and masked elements:} elements past $vl$ or masked off are left unchanged.
  \item \textbf{Strides:} a negative stride reads an array in reverse; a zero stride repeats one
        element.
  \item \textbf{Indexed and segment accesses:} gather/scatter follow the offset vector; segment
        loads split $(x, y)$ pairs and segment stores join them back.
  \item \textbf{Element widths and alignment:} byte elements, an element width different from
        SEW, and misaligned addresses all move the correct bytes.
  \item \textbf{Whole-register and mask accesses:} whole-register moves ignore $vl$; mask
        loads/stores move $\lceil vl/8 \rceil$ bytes.
  \item \textbf{Scalar/vector ordering:} a vector load sees an earlier scalar store, and vice
        versa, also on the pipelines.
  \item \textbf{Faults:} an invalid address stops the run with an access fault; a fault-only-first
        load shortens $vl$ instead; a misaligned register group is rejected as illegal.
\end{itemize}

\end{document}
